% Independent analytic insertion, 2026-10-08. No old source modified.
\subsection{A compiled finite prefix tree works for every fixed prime cutoff}
\label{sec:lower-k-uniform-product-interface-review}

The cutoff of the primes in the \emph{outer product sequence} is unrelated
to the first-prime cutoff in the \emph{modulus tree}. In particular the
outer cutoff need not exceed $1/12$, $1/50$, or the modulus presieve
exponent. We supply the finite-compilation argument needed to use the
new fifth/sixth density in the positive product theorem. This is a
local derivation from the named distribution theorems, not a numerical
replay or an independent proof of those external theorems.

\paragraph{The assertion with its quantifiers.}
Fix $a>0$ and a finite collection of labelled prime-product families,
with fixed bounded nonnegative label coefficients, at most a fixed
$k_*$ prime slots, every prime at least $N^a$, total product at most $N$,
and finitely many fixed monomial and ordering constraints. Let $w_N(m)$
be the resulting multiplicity and $X_N=\sum_m w_N(m)$. Suppose
$X_N=(J+o(1))N/\log N$. Let $K_{\rm end}$ be the limiting modulus-density
coefficient of the legal fourth/fifth/sixth tree, including its stated
order of replacements. For every \emph{strictly larger} $K>K_{\rm end}$,
\begin{equation}\label{eq:lower-k-uniform-product-conclusion}
 \sum_{\substack{m<N\\N-m\ {\rm prime}}}w_N(m)
 \le (KJ+o(1))\frac{C_NN}{\log^2N}.
\end{equation}
The error is uniform in the even integer $N$; its constants may depend
on $a$, the finite domains, and the chosen finite tree. The conclusion
also holds for a finite positive sum of such families. It does not
permit arbitrary bounded arithmetic reweightings, a cutoff tending to
zero with $N$, or an infinite adaptive tree. The assertion for $a\ge1$
is trivial here, so in the proof we take $0<a<1$.

\paragraph{Compile the new support before treating any outer primes.}
Choose a finite buffered version of the endpoint fourth/fifth/sixth
strategy whose limiting density coefficient is less than $K$ by a fixed
amount. This is possible because the numerical coefficient is used
with strict rounding slack, and the endpoint construction is a limit
of finite buffered trees. Keep all inherited coarse labels, the chosen
subset, every earlier membership decision, and the frozen level as
state. All these choices are independent of the actual outer primes
and of their progression errors. On a new selected branch at rank
$r=5$ or $6$, require, with a fixed positive margin,
\[
 \max\{1/2,W_r,G_I\}<\lambda<\vartheta<v(s_I),
 \qquad \vartheta<3/5.
\]
Here $\lambda$ is the level used in the odd cube tests and
$\vartheta$ is the slightly larger factorization parameter. The old
branches keep their own proved buffered factorization data.

At rank five the new test is imposed at that odd rank; the former
fifth test is replaced. At rank six membership is preserved, and only
later odd tests change. Thus every first exit remains odd. Write
$\lambda_{\mathcal T}(d)=\mu(d)\mathbf1_{d\in\mathcal T}$ for the
full raw support. The odd-first-exit identity gives a pointwise upper
sieve. This is not a claim that finite signed weights are monotone
under arbitrary support enlargement.

For a descendant $j>r$, all selected slots are among the first $r$,
so its remaining-selected square expression is exactly
$S_{j-1}+2t_j$. At an odd admitted rank the own cube implies this
square; at an even rank it follows from
\[
 S_{j-1}+2t_j\le S_{j-2}+3t_{j-1}.
\]
The first $r$ full and selected squares were checked in the displayed
initial gate. In the rank-six case the old fifth cube also proves that
the rank-six node was already present. Hence the complete descendant
has both full well support and every selected-remaining square gate
at the frozen level. This verifies the actual hypotheses of the
selected-subproduct factorization lemma; a bound on $s_I$ alone
would not do so. Subsequent primes are smaller than the prefix's last
prime, so the old first-three-slot data do not change.

Choose a presieve exponent $\omega>0$ below all active prefix cutoffs,
and put $w=N^\omega$. Rough support has bounded rank $M=O(1/\omega)$.
Refine the fixed coarse labels into fine prime intervals. At even
rough rank retain an outer box test on all required odd cubes; at odd
rough rank retain an inner box test. All coarse decisions are retained
exactly. The resulting coefficient $c(d)$ obeys
$c(d)\ge\lambda_{\mathcal T}(d)$ on squarefree rough $d$.
Discrepancies occur only near the finitely many frozen cube surfaces.
The tagged-prime estimate from the variable-prefix construction gives,
with fine width $\tau$,
\[
 V_N(w)\sum_d(c(d)-\lambda_{\mathcal T}(d))g_N(d)
 \ll V_N(z)C_{\mathcal T}\tau\omega^{-5}+o_N(V_N(z)).
\]
Changing at rank five or six only enlarges the finite list of frozen
levels in this argument. First deactivate upgrades on narrow coarse decision strips, reverting
there to the inherited legal rule; do not delete arbitrary moduli at an
even rank. Then apply this estimate within the fixed labels. No
discontinuous cap is interpolated across a strip.

Attach the small-prime fundamental-lemma weight $\beta^+$ to positive
rough coefficients and $\beta^-$ to negative ones. The signed
termwise inequalities for the two beta polynomials prove
\begin{equation}\label{eq:lower-k-compiled-majorant}
 \sum_{d\mid n}\rho_N(d)\ge\mathbf1_{(n,P_N(z))=1},
 \qquad |\rho_N(d)|\le1.
\end{equation}
The latter bound follows from the unique small/rough decomposition.
Their density error is $O(e^{-c/\eta}\omega^{-2})$ relative to $V_N(z)$.
The positive first-exit formula makes the rough density no larger
than the full raw density. Choosing the finite tree, beta precision,
and then fine width sufficiently carefully gives
\begin{equation}\label{eq:lower-k-compiled-density}
 \sum_d\rho_N(d)g_N(d)\le(K+o(1))\frac{C_N}{\log N}.
\end{equation}
This density computation contains no outer-product cutoff $a$.

\paragraph{Exact atoms, rather than factorability of an arbitrary weight.}
For each ordered vector of fine labels $J_1,\ldots,J_m$, the coefficient
is a fixed signed scalar times
\[
 \mathbf1_{\rm squarefree}
 (\chi_{J_1}*\cdots*\chi_{J_m})*\beta^\pm.
\]
The scalar contains the reciprocal repeated-label factorial exactly
once. There are finitely many atoms at fixed precision, and their
absolute scalar sum is bounded independently of $N$. The parity of
$\beta$ is fixed by the rough rank. This decomposition follows directly
from the boxed coefficient just constructed; it is not inferred from
pointwise membership of moduli in a well-factorable set.

For each required Type II scale, factor the upper endpoint vector of
one such box. The full and selected gates just proved give a fixed
partition of its prime slots into three groups. If group $i$ has $m_i$
slots, divide its prime convolution by $m_i!$; its coefficient is then
bounded by one, since there are at most $m_i!$ representations of an
integer. Restore the product of these group factorials in the atom's
scalar, without removing or reinserting the original repeated-label
factorial. Attaching $\beta^\pm$ to one rough group preserves the
one-bound by uniqueness of the small/rough factorization.

The factorization assignment may depend on the dyadic Type II scale,
but is fixed for every prime tuple in this box. This causes at most
logarithmically many distribution applications; it does not change
$\rho_N$. Removing rough square collisions costs
$(N+N^{\theta_*})N^{-\omega}(\log N)^C$ with a fixed $\theta_*<3/5$.
All complete moduli are below $N^{\theta_*}$. At rough rank at most
two their support is below $N^{2\rho+o(1)}$, where
$\rho=19/96$ and $2\rho<1/2$; classical/generalized BV suffices there.

\paragraph{Parameters can be chosen for arbitrarily small fixed $a$.}
After fixing the buffered tree with gate margin $e>0$, choose
\[
 0<\zeta\ll\min(e,a),\quad 12\zeta<a/2,\qquad
 \epsilon_{\rm disp}\ll\zeta,\quad \kappa\ll\epsilon_{\rm disp}.
\]
Next choose the beta level exponent $\ell$ so small that
$\ell\ll\epsilon_{\rm disp}$ and its density error is within the
rounding budget. With its resulting $\omega$ and $M$, choose $\tau$
so that $M\tau\ll\epsilon_{\rm disp}$ and the density-box error is
within that budget. Also require
$\ell+O(M\tau)\ll3/5-153/256$. Every choice is fixed before
$N\to\infty$. Shrinking $a$ changes these auxiliary choices and the
fixed number of prime slots, but changes neither the raw endpoint
tree nor its coefficient $K_{\rm end}$.

\paragraph{Distribution of the outer product: the individual rough slot.}
Take one positive smooth outer weight $W$ with every prime at least
$N^a/2$, and retain its exact coprime mass $X_W^{(d)}$.
Discard products below $N^{1-\kappa}$, with a power-saving error, and
use local product scales $N^{1-\kappa}/O(1)\le x\ll N$.
Apply the three-term Heath--Brown identity separately at each original
prime scale $R_i$. Discard every dyadic HB box which fails to meet
its original group interval before separating variables. The product
of its at most six dyadic lower endpoints is then
$\gg R_i/2^6\gg N^a$. This remains true on the separated rectangle.

Restrict the expanded variables to primes exceeding
$z_0=x^{1/(\log\log x)^3}$. This restriction is exact for the outer
prime slots once $N$ is large. Every nonunit HB variable is therefore
$z_0$-rough. If an individual variable lies in
$[x^{2\zeta},x^{1/3+\zeta/2}]$, its coefficient is rough $1$ or
rough $\mu$, possibly with one logarithm, not a grouped convolution.
The rough-SW lemma supplies its auxiliary-coprimality condition:
\[
 \left|\sum_{\substack{n\equiv r\ (q)\\(n,D)=1}}\alpha(n)
 -\frac1{\varphi(q)}\sum_{(n,Dq)=1}\alpha(n)\right|
 \ll_{A,a,\zeta}\frac{U\tau(D)}{\log^A x}.
\]
Its proof first sandwiches the rough indicator by fundamental-lemma
weights; the absolute difference controls the signed $\mu$ case.
For arbitrary auxiliary $D$, removing rough primes dividing $D$ costs
at most $(U/z_0+1)\omega(D)$, which is smaller than the displayed
bound for every fixed $A$. No part of this argument refers to the
modulus tree's fifth or sixth support.

The complementary HB factor has fixed divisor order depending on
$k_*$. Its large-coefficient tail is bounded after the modulus sum by
\[
 (\log N)^{-B}\sum_{m\ll N,\ m\ne N}
       \tau_C(m)\tau_C(|N-m|)
 \ll N(\log N)^{C'-B}.
\]
This uses Cauchy--Schwarz and divisor second moments. The isolated
$m=N$ term is $N^{\theta_*+o(1)}$; the mean tail uses ordinary divisor
moments and $\sum_{d\le N^{\theta_*}}\tau_C(d)/\varphi(d)
\ll\log^{C''}N$. Thus normalization reduces to bounded coefficients
without requiring rough SW for a new convolution.

\paragraph{Moving residue and the exceptional long cases.}
In L23 Proposition 5.2 use $a_{\rm res}=N$ and $A_0=Ux^{-\zeta}$.
The factorization starts with savings $3\zeta,3\zeta,2\zeta$;
replacing $A_0$ by $U$ costs $2\zeta,\zeta,(1+2\sigma)\zeta$.
For $\sigma=7/32$ the smallest surviving saving is $9\zeta/16$.
The small factor, fine endpoints and local coordinate change fit
within this slack. In the last target monomial the exponent of the
second modulus length is $5-2\sigma>0$, so dyadic subdivision does
not reverse an endpoint inequality. On each local scale
$\log N/\log x=1+O(\kappa)$ and $N<x^{1+\epsilon_{\rm disp}}$.
The source proposition has exactly this moving-residue supremum.

If there is no medium individual slot, all HB Mobius slots are short.
Every original group must contain a long slot; otherwise its product
is at most $x^{12\zeta+o(1)}<N^a/O(1)$. At most two long slots can
occur globally. Hence there are at most two original prime groups.
For two groups, there is one long slot in each and the remaining
product is at most $x^{20\zeta+o(1)}$. L23 Proposition 3.2 applies,
since $N_1,N_2>x^{3\epsilon_{\rm disp}}$ and
\[
 x^{3/5}<(x/x^{20\zeta+o(1)})^{2/3-3\epsilon_{\rm disp}}.
\]
For one group, its one-long and two-long alternatives use respectively
L23 Lemma 3.4 and Proposition 3.2. The one-long unrestricted mean is
replaced by the exact coprime mean using the bound
$\varphi(q)^{-1}\sum_{p\mid q,p>z_0}(t/p+1)
\ll_A t/(q\log^A x)$, already explicitly proved in the main text.
The residual exponent $20\zeta$ in this exceptional alternative is
independent of $k_*$ because the alternative itself forces $k\le2$
for that term; divisor orders and harmless constants elsewhere may
depend on the fixed total number of groups.

Mellin separation of $W$ gives only polylogarithmic twists of each
individual slot. Partial summation applies the interval-uniform rough
SW bound. In the two-long case use the arbitrary-subinterval version
in Maynard's proof of Proposition 5.2, not an arbitrary-coefficient
extension of that result. Mellin tails and prime powers use the same
divisor sum above. The finitely many atoms, prime patterns and
polylogarithmically many dyadic boxes are absorbed by increasing $A$.
We obtain the substantive bridge
\begin{equation}\label{eq:lower-k-smooth-exact-coprime}
 \sum_{(d,N)=1}\rho_N(d)
 \left\{\sum_{\prod p_i\equiv N\ (d)}W(\boldsymbol p)
           -\frac{X_W^{(d)}}{\varphi(d)}\right\}
 \ll_{A,a,W,\mathcal T}N(\log N)^{-A}.
\end{equation}

\paragraph{Mass, positivity and the same coefficient.}
Since $|\rho_N|\le1$ and the moduli are squarefree,
\[
 \sum_{d:p\mid d}\frac{|\rho_N(d)|}{\varphi(d)}
 \ll\frac{\log N}{p-1}.
\]
Consequently replacing $X_W^{(d)}$ by $X_W$ in
\eqref{eq:lower-k-smooth-exact-coprime} costs
$O_{a,k_*}(X_WN^{-a}\log N)$, with no modulus-level factor.
Apply \eqref{eq:lower-k-compiled-majorant} and then
\eqref{eq:lower-k-compiled-density}; the coefficient multiplying the
outer mass is precisely $K C_N/\log N$, not a rescaling by the
residual level of a switched node.

For sharp domains first use one nonnegative outer majorant in simplex
coordinates and $y=m/N$. Nonconstant monomial walls are simplex-null;
critical pure-product walls must instead be smoothed in $y$.
The radial mass lemma gives excess $\tau N/\log N$ and uniformly
bounded smooth seminorms at each fixed $\tau$. Apply the preceding
positive count theorem to this whole majorant before discarding its
excess mass. Small excess mass alone is not a bound for a signed
progression discrepancy. Let $N\to\infty$ first, and only afterwards
remove smoothing, the finite-tree approximation and its margins.
This proves \eqref{eq:lower-k-uniform-product-conclusion}.

\paragraph{A concrete boundary of the assertion.}
It would be false to replace the domain hypothesis by arbitrary bounded
outer arithmetic weights in the discrepancy statement. For example,
with $3\nmid N$ and one prime slot, retain only $p\equiv N\pmod3$.
At modulus $3$ its exact-coprime discrepancy equals $X/2$, where
$X\asymp N/(2\log N)$ on a fixed macroscopic interval. It has no
arbitrary logarithmic saving. This family fails the required tuple-domain
and SW hypotheses. The present fifth/sixth changes occur on the modulus
side and never make this prohibited change on the outer prime side.

\paragraph{Audit conclusion and its limits.}
The extension from the old modulus tree to the new finite fifth/sixth
rules is closed by the explicit compilation, descendant-square check
and parameter order above. In particular there is no additional rough
SW assumption and no use of a fixed-residue theorem at a growing
residue. The fixed-cutoff theorem follows from the same external
L23/Maynard distribution inputs and ordinary sieve/PNT inputs already
named in the manuscript. Their spectral proofs are not reproved here.
The finite-tree density certificates and ordinary density equality
are separately required and were not numerically replayed here.

The inspected recalibration fragment writes $K_1=1691/250=6.764$;
the current intended application uses $67633/10000=6.7633$.
Both are covered by the preceding argument if they exceed the
separately certified endpoint density. Their recalibration arithmetic
must use their own difference from $6.81$; the older $23/500$ must not
be silently reused for the newer value. This observation does not
alter either old source or the newer numerical work.
